\section{The learning flow under a discrimination field}
\label{app:flows}

\begin{figure}[h]
  \centering
  \includegraphics[width=\linewidth]{figures/learning_flows}
  \caption{\textbf{The discrimination field bends the flow of learning.}
  Arrows are the change $(\Fw, \Fmu)$ that one interaction induces in the
  receiver's fields; the background is $|(\Fw,\Fmu)|$ and the thin contours are
  the evidence $Z$. Centre: without discrimination the flow escapes the two
  dissonant quadrants into two symmetric consonant basins, separated by
  $\hmu=\hw$ (green). Left and right: a discriminating receiver evaluates the
  flow at $\hw+D$, which moves the separatrix and enlarges one basin at the
  other's expense --- ``distrust and disagree'' for an out-group emitter
  ($D=-d$), ``trust and agree'' for an in-group one ($D=+d$).}
  \label{fig:flows}
\end{figure}

\section{Derivation of the learning algorithm}
\label{app:derivation}

The receiver's belief before an interaction is the product of Gaussians
\eqref{eq:belief}. It observes $\sigma_e$, the emitter's opinion on issue
$\hat{x}$, and treats it as a report of an unknown true label $\sigma_T$
transmitted through a channel that flips the label with probability
$\varepsilon(z)$. Marginalising the joint distribution of labels and of the issue
as the emitter may have parsed it, $y_t$,
%
\begin{equation}
  L(\sigma_e\mid \hat{x}, u) \;=\;
  \sum_{\sigma_T}\int P(\sigma_e\mid \sigma_T, y_t, \hat{x}, u)\,
  P(\sigma_T\mid y_t, \hat{x}, u)\,P(y_t\mid \hat{x}, u)\,\mathrm{d}y_t ,
\end{equation}
%
with $P(\sigma_e\mid\sigma_T,z)=(1-\varepsilon(z))\delta_{\sigma_e,\sigma_T} +
\varepsilon(z)\delta_{\sigma_e,-\sigma_T}$ and
$P(\sigma_T\mid y_t,w)=\Theta(\sigma_T\, y_t\cdot w)$. Two modelling choices keep
the algebra Gaussian. Taking $\varepsilon(z)=\Phi(z)$ with $z$ Gaussian of mean
$\mu_{e|r}$ and variance $V_{e|r}$ makes the flip probability a bounded function
of a Gaussian variable while preserving conjugacy of the affective sector --- a
Beta distribution for $\varepsilon$, the tempting alternative, does not. Taking
$P(y_t\mid \hat{x},u)$ Gaussian with variance $v_p=\|w\|^{-2}$ lets an
inexperienced agent, whose weights are small, be correspondingly unsure that it
has parsed the issue as the emitter did. Together these give the likelihood
%
\begin{equation}
  L(\sigma_e\mid z, \hat{x}, w) \;=\;
  \Phi(z)\,\Phi(-w\cdot\hat{x}\sigma_e)
  + \bigl(1-\Phi(z)\bigr)\,\Phi(w\cdot\hat{x}\sigma_e),
  \label{eq:likelihood}
\end{equation}
%
and the covariance stays block diagonal in the two sectors if it starts so.

Averaging the likelihood over the current belief gives the evidence
%
\begin{equation}
  Z(\sigma_e\mid \hat{x}, D_t) \;=\; \mathbb{E}_{Q_t}\bigl[L(\sigma_e\mid z, u, \hat{x})\bigr]
  \;=\;\sum_{\sigma_T}\mathbb{E}_z\bigl[P(\sigma_e\mid\sigma_T,\hat{x},z)\bigr]\,
  \mathbb{E}_w\bigl[\Phi(\sigma_T\, w\cdot\hat{x})\bigr],
\end{equation}
%
and the Gaussian integrals collapse onto the two scaled fields
\eqref{eq:fields}, giving \eqref{eq:evidence}: the evidence depends on the
receiver's state only through $(\hw,\hmu)$. This is the symmetry that makes the
two sectors mirror images.

The exact Bayesian posterior is not Gaussian. Replacing it by the Gaussian of
maximum entropy relative to the prior, subject to matching the first two moments
of the posterior, gives updates in terms of derivatives of the log evidence,
%
\begin{equation}
  \hat{w}_{t+1} = \hat{w}_t + C_t\frac{\partial \ln Z}{\partial \hat{w}},
  \qquad
  C_{t+1} = C_t + C_t \frac{\partial^{2}\ln Z}{\partial\hat{w}\,\partial\hat{w}^{\!\top}} C_t,
\end{equation}
%
and identically in the affective sector with $\mu_{e|r}, V_{e|r}$ in place of
$\hat{w}, C$. Expressing the derivatives in the scaled fields yields
\eqref{eq:update-ideo}--\eqref{eq:update-aff} with the modulation functions
\eqref{eq:modulation}. The first is a Hebbian term $\sigma_e\hat{x}$ with a
tensorial, self-annealing rate; the rest is the online estimation of the channel.

\section{Genealogy of the learning rule}
\label{app:genealogy}

The update \eqref{eq:update-ideo}--\eqref{eq:update-aff} descends from the
variational optimisation of modulated Hebbian on-line learning
\citep{kinouchi1992optimal,kinzel1990improving}, given a Bayesian reading by
\citet{opper1996online,opper1998bayesian} and applied across architectures by
\citet{solla1999optimal}; the entropic-dynamics formulation we use is
\citet{caticha2020entropic}. Learning through a noisy channel whose noise level
must itself be inferred was analysed by \citet{kinouchi1993lower} and
\citet{biehl1995learning}; \citet{alves2016distrust} identified that inferred
level with distrust and made it dynamical, with applications to voting
\citep{caticha2015whom} and to panels of appellate judges
\citep{caticha2019trust}, and a mean-field treatment by \citet{simoes2018mean}.
The blame-attribution behaviour resembles the least-action rule of
\citet{mitchison1989bounds}, which suffers retarded learning
\citep{kabashima1994perfect}, whereas the optimised algorithm attains
generalisation error decaying as $\alpha^{-1}$ even under output noise
\citep{simonetti1996online}.

\section{The modulation functions over the plane}
\label{app:modulation}

\begin{figure}[h]
  \centering
  \includegraphics[width=0.8\linewidth]{figures/modulation_contours}
  \caption{$\Fmu, \Fw$ (top) and $\FV, \FC$ (bottom) over the plane of the opinion
  field $\hw$ and the distrust field $\hmu$. The green line is $\hmu=\hw$, and
  reflection across it exchanges the two sectors, $\Fw(x,y)=\Fmu(y,x)$ and
  $\FC(x,y)=\FV(y,x)$. Changes concentrate in the two dissonant quadrants ---
  agree-and-distrust at upper right, disagree-and-trust at lower left --- and vanish
  where the message is unsurprising.}
  \label{fig:modulation-contours}
\end{figure}

\section{The six discrimination cases}
\label{app:cases}

The matrix $D$ in \eqref{eq:disc-field} admits six inequivalent fillings, where
$+d$ is the tolerant entry and $-d$ the intolerant one (\S\ref{sec:discrimination}
and Appendix~\ref{app:sign}). Rows index the receiver's class, columns the
emitter's.

\newcommand{\dmat}[4]{\left(\begin{smallmatrix}#1&#2\\#3&#4\end{smallmatrix}\right)}

\begin{table}[h]
\centering\small
\begin{tabular}{cccc}
\toprule
case & $D/d$ & who discriminates & how \\
\midrule
1 & $\dmat{+1}{0}{0}{0}$   & $A$ only & in-group favouritism \\
2 & $\dmat{0}{-1}{0}{0}$   & $A$ only & out-group hostility \\
3 & $\dmat{+1}{-1}{0}{0}$  & $A$ only & both \\
4 & $\dmat{+1}{0}{0}{+1}$  & both     & in-group favouritism \\
5 & $\dmat{0}{-1}{-1}{0}$  & both     & out-group hostility \\
6 & $\dmat{+1}{-1}{-1}{+1}$ & both    & both \\
\bottomrule
\end{tabular}
\end{table}

All phase diagrams in the main text use case 6, in which the two classes play
symmetric roles.

\section{The sign of the discrimination field}
\label{app:sign}

The sign of $D$ is fixed by the algorithm and not free. By \eqref{eq:modulation},
$\Fmu=(1-2\Phi(\hw))g(\hmu)/Z$ is negative when $\hw>0$, so perceived agreement
drives distrust $\mu$ \emph{down}. A receiver that adds $+d$ to $\hw$ therefore
becomes more trusting (tolerant) and one that adds $-d$ becomes more distrustful
(intolerant). Case 6 with $d>0$ consequently gives in-group trust and out-group
distrust, i.e. the discriminatory phase, which is the convention used throughout.

The companion manuscript \citep{caticha2026discrimination} tabulates the in-group
entry as $-d$ while its text and figures place the discriminatory phases at
$d>0$; those two statements are inconsistent, and taken literally the tabulated
signs put the discriminatory phases at $d<0$, mirroring every map in $d$. The same
inconsistency appears in its description of the learning flow, where the panel
labelled $D_{e|r}=-d$ is called the more tolerant one although $-d$ is the
intolerant entry, and where the ``distrust and disagree'' basin is said to grow
for $d>0$ although adding $+d$ shrinks it. A single global sign reconciles
everything; we adopt it in the form that leaves \eqref{eq:disc-field} untouched.
The accompanying code implements both readings behind a flag, and
Figure~\ref{fig:sign} runs the same sweep under each: the two rows are mirror
images in $d$, and only the first matches the maps as published.

\begin{figure}[h]
  \centering
  \includegraphics[width=0.78\linewidth]{figures/sign_convention_comparison}
  \caption{\textbf{The two readings of the discrimination field.} Trust--class
  correlation and affective balance over the $(d,\fd)$ plane under the convention
  used in this paper (top) and under the companion manuscript's table read
  literally with its Eq.~25 (bottom). The physics is identical; which half of the
  $d$ axis is discriminatory is not.}
  \label{fig:sign}
\end{figure}

\section{Conventions in the order parameters}
\label{app:conventions}

Two choices in \eqref{eq:correlations} deserve statement.

\textbf{The class indicator.} We use $G_{IJ}=\kappa_I\kappa_J\in\{-1,+1\}$. A
$\{0,1\}$ indicator, which counts only same-class pairs, cannot represent
anti-correlation: under perfect reverse discrimination it returns
$-(N/2-1)/(N-1)$ rather than $-1$, and its range depends on $N$. Since the
reverse-discrimination half of the phase diagram is a substantive part of the
result, the signed indicator is the appropriate one. Both are implemented.

\textbf{Normalisation.} $\Rwmu$ and $\Rmuc$ sum $\eta_{I|J}+\eta_{J|I}$ over
unordered pairs, so dividing by $N(N-1)$ places them in $[-1,1]$. $\Rcw$ has a
single term per pair and the same divisor would cap it at $1/2$; we use
$2/(N(N-1))$ so that all three share the range $[-1,1]$ and the three colour
channels of Figure~\ref{fig:phase} are commensurate.

\section{Closed form for the balance aggregates}
\label{app:closed-form}

Enumerating triples costs $O(N^{3})$ per society, which is prohibitive across a
phase diagram. For any matrix $M$ with unit diagonal --- either $\rho$ or $\eta$ ---
%
\begin{equation}
  \sum_{\substack{(I,J,K)\\ \text{distinct, ordered}}} M_{IJ}M_{JK}M_{KI}
  \;=\; \operatorname{tr}(M^{3}) - 3\Bigl(\textstyle\sum_{I,J}M_{IJ}M_{JI} - N\Bigr) - N ,
\end{equation}
%
and each unordered triple appears six times, three as the forward cycle and three
as the reverse, so
$\langle b\rangle = \bigl[\text{right-hand side}\bigr]/\bigl(6\binom{N}{3}\bigr)$.
For symmetric $M$ this is $\BI$; for the asymmetric trust matrix the two cycle
orientations are exactly the two terms of $b^{A}$, so the same expression gives
$\BA$. Both are batched matrix products. The identity is verified against
enumeration in the test suite.

\section{Parameters}
\label{app:parameters}

The model needs six numbers. One is recovered from the companion manuscript, one
is calibrated against its published figures, two are prior conventions, and two
are free choices; all six are listed with their provenance rather than left
implicit, since the manuscript states none of them. The accompanying code exposes
each of them and its calibration script reproduces the comparisons below.

\begin{table}[h]
\centering\small
\begin{tabular}{llll}
\toprule
symbol & meaning & value & provenance \\
\midrule
$K$ & dimension of the issue embedding & $30$ & recovered (see below) \\
$N$ & agents in the society & $40$ & chosen; converged (see below) \\
$P$ & issues on the agenda & $5$ / $100$ & chosen: $\alpha=P/K$ either side of $1$ \\
$\Delta t$ & interactions per ordered pair & $500$ & calibrated (see below) \\
$C_0, V_0$ & initial uncertainties & $I$, $1$ & uninformative prior \\
$\mu_0$ & initial distrust & $\mathcal{U}(-1,1)$ & gives $\BI\approx\BA\approx0$ at $t=0$ \\
\bottomrule
\end{tabular}
\end{table}

\paragraph{$K$ is recovered, not chosen.} The companion manuscript reports
balance trajectories at
$\alpha \in \{0.03, 0.17, 0.23, 0.33, 0.50, 0.67, 1.67, 3.33, 333.33\}$. These are
exactly $P/30$ for $P\in\{1,5,7,10,15,20,50,100,10^{4}\}$, which fixes $K=30$;
the language-agent protocol's thirty issues agrees.

\paragraph{$\Delta t$ is calibrated against the published trajectories.} The
dynamics anneals rather than reaching a stationary state, so the measurement time
matters and the companion manuscript does not state it (its text carries a
literal placeholder in place of the value). We scan $\Delta t$ and compare the
endpoints of all nine balance trajectories against the published ones:

\begin{table}[h]
\centering\small
\begin{tabular}{lccccc}
\toprule
$\Delta t$ (interactions per ordered pair) & 60 & 125 & 250 & \textbf{500} & 1000 \\
rms distance to published endpoints & 0.63 & 0.31 & 0.15 & \textbf{0.13} & 0.15 \\
\bottomrule
\end{tabular}
\end{table}

The minimum is flat-bottomed and sits at $\Delta t=500$, which we use throughout.
Seven of the nine endpoints then agree to within about $0.1$, and $\BI$ matches
the three largest agendas almost exactly, but no single $\Delta t$ reproduces all
nine: two middle curves ($\alpha=0.23$, $0.33$) sit above ours, and the two
largest agendas have a lower published $\BA$ than we obtain at any $\Delta t$ that
fits the rest. Either the society size differs from ours or those middle curves,
which overlap within a few pixels in the printed figure, cannot be digitised
reliably; we cannot distinguish the two. Dropping the ambiguous readings does not
move the selected value. All the $\Delta t$-independent checks pass: the
above/below-diagonal ordering with its crossover at $\alpha\approx1.7$, the sharp
sign flip of $\Rmuc$ at $d\approx0$, and the $\Rcw$ wedge confined to large $d$
and large $\fd$.

\paragraph{$N=40$ is large enough.} Repeating five representative points of the
plane at $N\in\{20,30,40,60\}$ with twelve realisations each, the order parameters
at $N=60$ differ from those at $N=40$ by at most $0.02$ ($\Rmuc$, $\Rcw$), $0.03$
($\BI$) and $0.05$ ($\BA$), while $N=20$ differs by up to $0.24$. The maps are
therefore converged in $N$ at $N=40$, and the cost of going further is steep: total
work grows as $N^{3}$.

\paragraph{Numerical safeguards.} Three floors are applied, none of which the
model prescribes: on the evidence $Z$, whose exact modulation functions diverge
in the deep dissonance corners; on the affective variance $V$, which $\FV<0$
drives monotonically towards zero; and on the covariance update, clipped to the
boundary at which $C+a(C\hat{x})(C\hat{x})^{\!\top}$ ceases to be positive
semi-definite. The last is the only one that can alter a trajectory, and it never
triggered in the runs reported here.

